\documentclass{article}
\usepackage[utf8]{inputenc}
\usepackage[T1]{fontenc}
\usepackage{amsmath, amssymb, amsthm}
\usepackage{geometry}
\usepackage{hyperref}
\usepackage{bookmark}
\usepackage{tikz}

\usepackage{titling}


\newcommand{\subtitle}[1]{
    \posttitle{\par\end{center}\begin{center}\large#1\end{center}\vskip0.5em}
}


\setlength{\parindent}{0pt}
\setlength{\parskip}{0.8em}

\title{Correction du CC3 2022}

\subtitle{MAP 151} 

\author{Axel Lavigne}
\date{\today}


\newcommand{\correction}[0]{
    \vspace{1mm}
    \textbf{Correction}
    \hrulefill
}

\newcommand{\smalltext}[1]{
    \scriptsize
    \textit{#1}
    \normalsize
}

\begin{document}

\maketitle


$\mathbf{Exercice}$1:

On considère la fonction $f:x\mapsto\frac1{1+x^2}.$ L'objectif est d'évaluer l'intégrale
$$I=\int_0^1\frac{1}{1+x^2}.$$


(1) Donner la valeur exacte de $I.$

\correction

\begin{align*}
    I &= \int_0^1\frac{1}{1+x^2}dx \\
    &= \left[ \arctan(x) \right]_0^1 \\
    &= \arctan(1) - \arctan(0) \\
    &= \frac{\pi}{4} - 0 \\
    &= \frac{\pi}{4}.
\end{align*}


\hrulefill


(2) Evaluer numériquement avec 4 sous-intervalles en utilisant:


a. la méthode des trapèzes. 

b. la méthode de Simpson. On pensera à donner la formule du cours.

On demande une précision à $10^{-5}$ près


\correction

Notons $x_i$ les points séparant les sous-intervalles. 
On a $x_0=0, x_1=\frac14, x_2=\frac12, x_3=\frac34$ et $x_4=1.$

On note $h=\frac14$ la longueur des sous-intervalles.

a. La méthode des trapèzes consiste à approximer l'intégrale par l'aire d'un trapèze.
On a la formule suivante:

\begin{equation*}
    \int_a^b f(x) \, dx \approx \frac{h}{2} \left[ f(a) + f(b) + 2 \sum_{i=1}^{n-1} f(x_i) \right].
\end{equation*}

\begin{equation*}
    \sum_{i=1}^{n-1} f(x_i) \approx 2,381176471
\end{equation*}

En appliquant cette formule, on obtient:

\begin{equation*}
    \frac14 \left[ 1 + 0.39024 + 2 \sum_{i=1}^{n-1} f(x_i) \right] \approx 0.782794118
\end{equation*}


b. La méthode de Simpson consiste à approximer l'intégrale par l'aire d'un polynôme de degré 2.

\begin{equation*}
    \int_a^b f(x) \, dx \approx \frac{h}{6} \left[ f(a) + f(b) + 2 \sum_{i=1}^{n-1} f(x_{i}) + 4 \sum_{i=1}^{n} f(x_{i-\frac12}) \right].
\end{equation*}

avec $x_{i-\frac12} = \frac{x_{i-1} + x_i}{2}$.

En faisaint les calculs, on obtient:

\begin{equation*}
    \int_a^b f(x) \, dx \approx \frac{h}{6} \left[ f(a) + f(b) + 2 \sum_{i=1}^{n-1} f(x_{i}) + 4 \sum_{i=1}^{n} f(x_{i-\frac12}) \right] = 0,782794118
\end{equation*}

Soit la meme valeur que la méthode des trapèzes.


\hrulefill

(3) 

a. Justifier rapidement que $f$ est trois fois dérivable sur l'intervalle [0;1] et donner
l'expression des dérivées $f^\prime,f^{\prime\prime}$ et $f^{(3)}.$

\correction

a. La fonction $f$ est trois fois dérivable sur $[0,1]$. car elle est la composée de fonctions dérivables.


Ses dérivées successives sont:

\begin{align*}
    f'(x) &= -\dfrac{2x}{(1 + x^2)^2} \\
    f''(x) &= \dfrac{6x^2 - 2}{(1 + x^2)^3} \\
    f^{(3)}(x) &= -\dfrac{24x(x^2 - 1)}{(1 + x^2)^4}
\end{align*}


\hrulefill

b. Donner la formule de l'erreur pour la méthode des trapèzes.

c. En déduire un rang à partir duquel l'erreur de quadrature est inférieure à
$10^{-6}.$


\correction

b. La formule de l'erreur pour la méthode des trapèzes est donnée par:

\begin{equation*}
    \left| \int_a^b f(x) \, dx - \frac{h}{2} \left[ f(a) + f(b) + 2 \sum_{i=1}^{n-1} f(x_i) \right] \right| \leq \frac{M_2}{12}h^2(b-a)
\end{equation*}

avec $M_2 = \max_{x \in [a,b]} |f''(x)|.$

c. On a $M_2 = \max_{x \in [0,1]} |f''(x)|$. 

Comme $f''(x)$ est croissante sur $[0,1]$, 

$M_2 = |f''(1)| = \dfrac{8}{4} = \dfrac{1}{2}$.

L'erreur de quadrature est majorée par:
\begin{equation*}
    \text{Erreur} \leq \dfrac{M_2}{12}h^2(b - a)
\end{equation*}

On veut que $\text{Erreur} \leq 10^{-6}$, donc:
\begin{align*}
    \dfrac{M_2}{12}h^2 &\leq 10^{-6} \\
    h^2 &\leq \dfrac{12 \times 10^{-6}}{M_2} \\
    h &\leq \sqrt{24 \times 10^{-6}} \approx 0.004898979486
\end{align*}

\hrulefill

$\mathbf{Exercice}$2:

On considère le problème de Cauchy suivant:

$$
\begin{cases}y'(t)+ty(t)=\left(t^2+1\right)e^{-\frac{t^2}{2}},\:\forall t\in\mathbb{R},\\y(0)=1.\end{cases}
$$
1. Résoudre le problème de Cauchy.

\correction

On cherche une solution de la forme $y(t) = y_h(t) + y_p(t)$, 
où $y_h(t)$ est la solution de l'équation homogène 
et $y_p(t)$ est une solution particulière de l'équation non homogène.

L'équation homogène associée est $y'(t) + ty(t) = 0$.

les solution sont triviales, on a $y_h(t) = \lambda e^{-\frac{t^2}{2}}$.

Pour trouver une solution particulière, on utilise la méthode de variation de la constante.

On pose $y_p(t) = \lambda(t) e^{-\frac{t^2}{2}}$.

il s'agit de trouver $\lambda(t)$ tel que $y_p(t)$ soit solution de l'équation non homogène.

On poses alors

\begin{align*}
    \lambda'(t) e^{-\frac{t^2}{2}} &= \left(t^2 + 1\right)e^{-\frac{t^2}{2}} \\
    \lambda'(t) &= \left(t^2 + 1\right)e^{-\frac{t^2}{2}} e^{\frac{t^2}{2}} \\
    \lambda'(t) &= \left(t^2 + 1\right)
\end{align*}

On intègre $\lambda'(t)$ pour trouver $\lambda(t)$:

\begin{align*}
    \int \lambda'(t) \, dt &= \int \left(t^2 + 1\right) \, dt \\
    \lambda(t) &= \frac{t^3}{3} + t + C
\end{align*}

On fixe $C = 0$.

On a donc $y_p(t) = \left(\frac{t^3}{3} + t\right) e^{-\frac{t^2}{2}}$.

Ainsi la solution générale est $y(t) = \lambda e^{-\frac{t^2}{2}} + \left(\frac{t^3}{3} + t\right) e^{-\frac{t^2}{2}}$.

Il reste à trouver $\lambda$ en utilisant la condition initiale $y(0) = 1$.

\begin{align*}
    y(0) &= \lambda e^{-\frac{0^2}{2}} + \left(\frac{0^3}{3} + 0\right) e^{-\frac{0^2}{2}} \\
    1 &= \lambda + 0 \\
    \lambda &= 1
\end{align*}

La solution du problème de Cauchy est donc 

\begin{equation*}
    y(t) = e^{-\frac{t^2}{2}} + \left(\frac{t^3}{3} + t\right) e^{-\frac{t^2}{2}}
\end{equation*}

\hrulefill

2. On se place sur l'intervalle [0;2].

Mettre en place la méthode d'Euler explicite avec un pas $h=\frac{1}{2}$.
On notera $y_n$ l'élément de rang $n$ de la suite utilisée.  On prendra soin
de donner tous les détails de la construction.

\correction

On a $h = \frac{1}{2}$ et $n = 4$.

Notons $y_0 = 1$, $t_0 = 0$.

On définit $t_{n+1} = t_n + h$.
ce qui est équivalent à $t_{n} = nh \rightarrow t_{n} = \frac{n}{2}$.

On a l'équation suivante:

\begin{align*}
    y'(t)+ty(t)&=\left(t^2+1\right)e^{-\frac{t^2}{2}} \\
    y'(t) &= \left(t^2+1\right)e^{-\frac{t^2}{2}} - ty(t)
\end{align*}

On vois apparaite la forme $$y'(t) = f(t,y(t))$$
dont le shema d'Euler explicite est donné par:
$$
\begin{cases}
    y_{n+1} = y_n + h f(t_n,y_n) \\
    t_n = nh
\end{cases}
$$


On peut donc en déduire le shema d'Euler explicite suivant :

\begin{equation*}
    \begin{cases}
        y_{n+1} = y_n + h \left[ \left(t_n^2 + 1\right)e^{-\frac{t_n^2}{2}} - t_n y_n \right] \\
        y_{0} = 1
    \end{cases}
\end{equation*}

On peut alors calculer les valeurs de $y_n$ pour $n = 0, 1, 2, 3, 4$.

On obtient les valeurs suivantes:
\begin{align*}
    y_0 &= 1 \\
    y_2 &= 1.5 \\
    y_3 &= 1.676560564 \\
    y_4 &= 1.444810942 \\
    y_5 &= 0.888762995
\end{align*}




\hrulefill

3. On se place sur l'intervalle [0;2]. Mettre en place la méthode d'Euler implicite avec
un pas $h=\frac12.$ On notera $x_n$ l'élément de rang $n$ de la suite utilisée.

\correction

de meme que pour la méthode explicite, on a $h = \frac{1}{2}$ et $n = 4$.

Nous pouvons alors dresser le shema d'Euler implicite suivant:

\begin{equation*}
    \begin{cases}
        y_{n+1} = y_n + h \left[ \left(t_{n+1}^2 + 1\right)e^{-\frac{t_{n+1}^2}{2}} - t_{n+1} y_{n+1} \right] \\
        y_{0} = 1
    \end{cases}
\end{equation*}

En reformulant l'équation pour avoir $u_{n+1}$ en fonction de $u_n$, on obtient:

\begin{equation*}
    y_{n+1} = \dfrac{y_n + h \left[ \left(t_{n+1}^2 + 1\right)e^{-\frac{t_{n+1}^2}{2}} \right]}{1 + h t_{n+1}}
\end{equation*}

On peut alors calculer les valeurs de $y_n$ pour $n = 0, 1, 2, 3, 4$.

On obtient les valeurs suivantes:

\begin{align*}
    y_0 &= 1 \\
    y_2 &= 1.241248451 \\
    y_3 &= 1.231852741 \\
    y_4 &= 1.005378857 \\
    y_5 &= 0.671858533
\end{align*}


\hrulefill

4. Comparez les valeurs obtenues aux questions précédentes (pour l'approximation de
$y(2))$ avec la valeur exacte $y(2)$ et commenter.

\correction

les valeurs exactes à chauqes étapes sont les suivantes:

\begin{align*}
    y(0) &= 1 \\
    y(0.5) &= 1.360516058 \\
    y(1) &= 1.415238206 \\
    y(1.5) &= 1.176865194 \\
    y(2) &= 0.766899938
\end{align*}

Nous pouvons alors comparer les valeurs obtenues avec les méthodes d'Euler explicite et implicite avec la valeur exacte.

Nous avons le tableau suivant:

\begin{center}
\begin{tabular}{|c|c|c|c|}
    \hline
    Méthode & $y(2)$ & Erreur \\
    \hline
    Exact & 0.766899938 & 0 \\
    Euler explicite & 0.888762995 & 0.121863057 \\
    Euler implicite & 0.671858533 & 0.095041405 \\
    \hline
\end{tabular}
\end{center}

L'érreur se calcule en faisant la différence entre la valeur exacte et la valeur obtenue.

On remarque que la méthode d'Euler explicite est plus précise que la méthode d'Euler implicite.

\hrulefill


\end{document}


